%% ma: L11-B15-0037
%% lop: 11
%% chuong: 5
%% bai: 15
%% dang: 11.15.3
%% loai: TN4
%% muc: VD
%% dap_an: "A"
%% nguon: "tai-lieu-hoc-tap-mon-toan-lop-11-theo-chuong-trinh-giao-duc-pho-thong-2018.pdf (Lớp Toán Thầy Duy, Phạm Lê Duy, THPT Chu Văn An), Chương 3 Bài 1, Phần 1, Câu 26 (tr.324)"
%% kien_thuc: "Lớp 11 Bài 15 (giới hạn của dãy số)"
%% trang_thai: cach-ly
%% ly_do: "Đề và đáp án gốc không khớp: theo đề, lim = 3/64 nên a·b = 192, không có trong bốn phương án; đáp án gốc A (1728) ứng với 27/64, tức tử là (3n-1)^3. Gợi ý sửa: đổi tử thành (3n-1)^3 (khi đó lim = 27/64, a·b = 1728) hoặc đổi một phương án thành 192."
%% ghi_chu: "Câu có lỗi đề/đáp án của tài liệu gốc."
\begin{ex}
Dãy số $(u_n)$ với $u_n=\dfrac{(3n-1)(3-n)^2}{(4n-5)^3}$ có giới hạn bằng phân số tối giản $\dfrac ab$. Tính $a\cdot b$.
\choice{\True $1728$}{$68$}{$32$}{$128$}
\loigiai{\sloppy Chia tử và mẫu cho $n^3$: \[u_n=\dfrac{\left(3-\frac1n\right)\left(\frac3n-1\right)^2}{\left(4-\frac5n\right)^3}\to\dfrac{3}{64},\] nên $a\cdot b=192$.}
\end{ex}
